% Article VII — first editorial draft, 10 September 2026.
% Opening integrated with the proofs of revision 06.
% Provenance and conditions: PLAN_RESIDENCIAS_DEMOSTRATIVAS.md.
% Internal references point to proofs included in the body.

\begin{center}\bfseries Abstract\end{center}
\begingroup
\fontsize{9.5}{10.7}\selectfont
\setlength{\parskip}{1.5pt plus .5pt minus .5pt}
\noindent
Torsional amplitude is determined as an explicit functional of the
material state through the spinorial action, elimination of contorsion,
and metric variation. A material sector is constructed in which this
amplitude is positive and conserved. Its energetic contribution produces
\(a^{-6}\) dilution and, in the flat homogeneous realization with
\(\Lambda_{\mathrm{eff}}=0\), positive densities, and \(w<1\), a unique
transverse bounce threshold. For dust, an exact solution is obtained for
all time, with finite curvature and Levi--Civita causal geodesics complete
in both directions.

The construction starts from a formal kernel called Triadic Modular
Holography, henceforth HMT. \emph{All Possible Paths}, henceforth APP,
comprises arithmetic paths on \((\mathbb Z/9\mathbb Z)^2\), with
adjacency given by the additive Cayley graph with cardinal generators
and a toroidal cellular realization, distinguishing sum and product
operations and evaluations; TRIT is the unit
of oriented signature \(\{-1,0,+1\}\), determining local regimes; the
\emph{Topological Phase Kernel}, henceforth TPK, composes the dynamics of
transport and memory. The enriched state preserves sheets, orientation,
remainder, carry, and history during phase return. Dimensional Mechanics
expresses its internal results through physical quantities and
observables, preserving the structural provenance of the couplings.

Gravitational normalization is first developed through the longitudinal
composition \(L=(5759/23040)\alpha^{16}L_*\) and its positive radial
metric. In the realization defined by these operations and by the
identification of the reduced gravitational radius, reciprocity with the
action section fixes \(G=c^3L^2/\hbar_{\rm ret}\) in every admissible
mode. The proof preserves the full archive and establishes the conditions
for static and dynamic reduction of memory. The dimensional choice
\(L_*=1\,\mathrm m\), the longitudinal and metric rules, and the
radial identification are stated in the theorem itself.

Route transport admits a Cartan realization compatible with composition,
inversion, and refinement. Two material actions are distinguished on that
realization: the minimizing extension of the memory screen and the
spinorial realization. Each retains its own variation with respect to
the connection. In the spinorial sector, affine in contorsion, inversion
of the Holst and torsion operators determines the effective correction;
for the minimal extension, the nonlinear dependence of the current is
preserved and a local criterion is proved through the Hessian. The
current, quartic observable, and its normalization explain the value of
the amplitude as well as its role in cosmological evolution.

Persistence of the bounce under perturbations is proved through distinct
local bounds, whereas causal completeness belongs to the explicit global
dust solution. Residual balances, horizon information, and the observer's
operations preserve the conditions on sources, couplings, and charts
that they use. The relational response of the joint state is connected
to advanced and retarded propagation through Green operators for finitely
supported sources, associated with the covariant difference constructed
from route transport and a boundary
compatibility condition. The result relates generated memory, its
geometric realization, and the dynamical and informational consequences
of that realization.

The composition of action, the Barbero--Immirzi area increment, and the
horizon determines the conversion
\(T_{\rm cos}(L_\alpha)/\Theta_{\rm clk}=\log3/(2\pi)\).
The angular response Hamiltonian yields a finite free-energy balance
through relative entropy and a first law at fixed radius. A relative
thermometric section also characterizes the positive transducer and
the criterion for identifying another reader on that same section.

\noindent\textbf{Keywords:} gravitation; torsion; arithmetic memory; field equations; Planck length; transition from contraction to expansion.
\par\endgroup
